\documentclass[11pt,a4paper]{article}
\usepackage[T1]{fontenc}
\usepackage[utf8]{inputenc}
\usepackage{amsmath,amssymb,amsthm}
\usepackage[margin=1in]{geometry}
\usepackage[colorlinks=true,linkcolor=blue,urlcolor=blue]{hyperref}
\theoremstyle{plain}
\newtheorem{theorem}{Theorem}
\newtheorem{lemma}[theorem]{Lemma}
\newtheorem{proposition}[theorem]{Proposition}
\newtheorem{corollary}[theorem]{Corollary}
\theoremstyle{definition}
\newtheorem{remark}[theorem]{Remark}

\title{The reduction of OEIS A196523 modulo $3$}
\author{Adrian Perez Fontelles\\ \small Independent researcher}
\date{31 August 2026}
\begin{document}
\maketitle

\begin{abstract}
OEIS A196523 is defined by the functional equation $A(x)=x+x\,A^{4}(x)$, where
$A^{4}$ is the $4$-th compositional iterate, and carries a conjecture of
P.~D.~Hanna that $a(n)\equiv1\pmod{3}$ for every $n\ge1$. It is true. The equation determines the coefficients one at a
time, with $a(N)$ entering with coefficient $1$, so it determines them modulo $3$ as
well. Modulo $3$ the series $x/(1-x)$ satisfies the equation exactly, because its
iterates are $x/(1-jx)$ and $4\equiv1$ modulo $3$. Uniqueness then identifies the
reduction. The argument is an identity between rational functions over
$\mathbb{Z}/3\mathbb{Z}$; nothing is truncated or estimated.
\end{abstract}

\noindent\small 2020 Mathematics Subject Classification. 11B50, 05A15, 13F25.\normalsize

\section{The sequence and the conjecture}

OEIS A196523 is ``G.f. A(x) satisfies A(x) = x + x*A(A(A(A(x)))).''. It has offset $1$ and begins
\[
1, 1, 4, 28, 262, 2944, 37666,\ \dots
\]
Writing $A(x)=\sum_{n\ge1}a(n)x^{n}$ and $A^{\,i}$ for the $i$-th compositional
iterate of $A$, the entry defines $A$ by
\begin{equation}\label{eq:fe}
A(x)\;=\;x+x\,A^{4}(x).
\end{equation}
This is the only input taken from the entry.

The entry separately carries the following comment:
\begin{quote}\small
Conjecture: a(n) == 1 (mod 3) for n >= 1. --- \emph{Paul D. Hanna}
\end{quote}
As of the ``Last modified'' line on the live entry (Jun 13 2026 09:32:51, revision 30) the statement is
still recorded as a conjecture, and nothing on the entry records it as settled either way.

\section{The equation determines the sequence}

\begin{lemma}\label{lem:jaN}
Let $A(x)=\sum_{n\ge1}a(n)x^{n}$ with $a(1)=1$. Then for every $i\ge1$ and every $N\ge1$,
\[
\bigl[x^{N}\bigr]A^{\,i}=i\,a(N)+\bigl(\text{a polynomial in }a(1),\dots,a(N-1)\bigr).
\]
\end{lemma}

\begin{proof}
Induction on $i$. For $i=1$ the statement is $\bigl[x^{N}\bigr]A=a(N)$. Suppose it holds
for $i$ and put $B=A^{\,i}$, so $\bigl[x^{1}\bigr]B=1$ and
$\bigl[x^{N}\bigr]B=i\,a(N)+\cdots$. Then
$A^{\,i+1}=A\circ B=\sum_{m\ge1}a(m)B^{m}$, so
\[
\bigl[x^{N}\bigr]A^{\,i+1}=\sum_{m=1}^{N}a(m)\,\bigl[x^{N}\bigr]B^{m} .
\]
The term $m=N$ contributes $a(N)\bigl(\bigl[x^{1}\bigr]B\bigr)^{N}=a(N)$; the term $m=1$
contributes $\bigl[x^{N}\bigr]B=i\,a(N)+\cdots$; and for $2\le m\le N-1$ every monomial of
$\bigl[x^{N}\bigr]B^{m}$ is a product of $m\ge2$ coefficients of $B$ whose indices sum to
$N$, so each index is at most $N-1$ and no $a(N)$ occurs. Adding up gives $(i+1)a(N)$ plus
terms in $a(1),\dots,a(N-1)$.
\end{proof}

\begin{lemma}\label{lem:unique}
There is exactly one sequence of integers $a(1),a(2),\dots$ satisfying \eqref{eq:fe}, and
for each $N\ge2$ the equation read at order $x^{N}$ has the form
$a(N)=(\text{a polynomial with integer coefficients in }a(1),\dots,a(N-1))$.
\end{lemma}

\begin{proof}
Comparing coefficients of $x$ in \eqref{eq:fe} gives $a(1)=1$, since $x\,A^{4}(x)$ has
no linear term. Let $N\ge2$. The left side contributes $a(N)$. On the right,
$\bigl[x^{N}\bigr]\bigl(x\,A^{4}\bigr)=\bigl[x^{N-1}\bigr]A^{4}$, which by
Lemma~\ref{lem:jaN} involves only $a(1),\dots,a(N-1)$. So the coefficient of $a(N)$ in the
order-$x^{N}$ equation is $1$, the equation solves for $a(N)$ over $\mathbb{Z}$, and
induction gives existence and uniqueness.
\end{proof}

\begin{corollary}\label{cor:mod}
Let $m\ge2$. If $\bar A\in(\mathbb{Z}/m\mathbb{Z})[[x]]$ satisfies $\bar A=x+O(x^{2})$ and
\eqref{eq:fe} with all coefficients read modulo $m$, then
$a(n)\equiv\bigl[x^{n}\bigr]\bar A\pmod m$ for every $n$.
\end{corollary}

\begin{proof}
Reduction modulo $m$ is a ring homomorphism, so it carries the integer recursion of
Lemma~\ref{lem:unique} to the same recursion over $\mathbb{Z}/m\mathbb{Z}$, in which the
coefficient of $a(N)$ is still $1$ and hence still invertible. That recursion has a unique
solution, and both $a\bmod m$ and the coefficients of $\bar A$ solve it.
\end{proof}

\section{The reduction modulo $3$}

Write $f(x)=\dfrac{x}{1-x}$ over $\mathbb{Z}/3\mathbb{Z}$.

\begin{lemma}\label{lem:iter}
$f^{\,i}(x)=\dfrac{x}{1-ix}$ for every $i\ge1$.
\end{lemma}

\begin{proof}
True for $i=1$. If $f^{\,i}(x)=x/(1-ix)$ then
\[
f^{\,i+1}(x)=f\bigl(f^{\,i}(x)\bigr)
=\frac{\dfrac{x}{1-ix}}{1-\dfrac{x}{1-ix}}
=\frac{x}{(1-ix)-x}=\frac{x}{1-(i+1)x} .
\]
\end{proof}

\begin{lemma}\label{lem:sat}
Over $\mathbb{Z}/3\mathbb{Z}$, $f$ satisfies $f(x)=x+x\,f^{4}(x)$.
\end{lemma}

\begin{proof}
By Lemma~\ref{lem:iter}, $f^{4}(x)=\dfrac{x}{1-4x}$. Since $4\equiv1$ modulo
$3$ --- this is the whole content of the choice of modulus, as $3=4-1$ --- we have
$1-4x=1-x$ in $(\mathbb{Z}/3\mathbb{Z})[x]$, so $f^{4}(x)=\dfrac{x}{1-x}$.
Therefore
\[
x+x\,f^{4}(x)=x+\frac{x^{2}}{1-x}
=\frac{x(1-x)+x^{2}}{1-x}=\frac{x}{1-x}=f(x).
\]
\end{proof}

\begin{theorem}\label{thm:main}
Over $\mathbb{Z}/3\mathbb{Z}$ the reduction of $A$ is $x/(1-x)$; equivalently
$a(n)\equiv1\pmod{3}$ for every $n\ge1$. This is the conjecture of Section~1.
\end{theorem}

\begin{proof}
By Lemma~\ref{lem:sat} the series $f=x/(1-x)$ satisfies the reduced equation, and
$f=x+O(x^{2})$. By Corollary~\ref{cor:mod} the reduction of $A$ equals $f$, every
coefficient of which is $1$.
\end{proof}

\begin{remark}
The proof used only $4\equiv1$ modulo $3$. The same argument shows that for every
$i\ge2$ the sequence defined by $A(x)=x+x\,A^{\,i}(x)$ has all of its terms congruent to
$1$ modulo $i-1$. It also shows where the argument stops: if $m\nmid i-1$ then $1-ix$ and
$1-x$ differ in $(\mathbb{Z}/m\mathbb{Z})[x]$, so $x/(1-x)$ is not a solution of the
reduced equation and this route says nothing modulo $m$.
\end{remark}

\section{Verification}

Three checks, each independent of the algebra above.

First, the recursion of Lemma~\ref{lem:unique} was run in exact rational arithmetic and
its output compared against the entry's DATA at every one of the $20$ published
indices; the values agree exactly, and every one is an integer. A recursion that reproduced
only some of them would mean the functional equation had been transcribed wrongly.

Second, at every step the coefficient of $a(N)$ was computed rather than assumed, by
evaluating the order-$x^{N}$ equation at $a(N)=0$ and at $a(N)=1$ and subtracting; it came
out $1$ at every $N$ up to $23$, as Lemma~\ref{lem:unique} requires. The run was continued
to $n=23$ and every term reduced modulo $3$; all residues are $1$.

Third, both sides of \eqref{eq:fe} were expanded independently as formal power series over
$\mathbb{Z}/3\mathbb{Z}$ to order $30$, starting from the coefficients of $x/(1-x)$
rather than from the closed form, and they agree at every order; the iterates $f^{\,i}$
were confirmed to have coefficients $i^{\,n-1}$ for $i\le6$.

\begin{thebibliography}{9}
\bibitem{oeis} The OEIS Foundation, \emph{The On-Line Encyclopedia of Integer
Sequences}, \url{https://oeis.org}, 2026. Sequence A196523.
\bibitem{stanley} R.~P.~Stanley, \emph{Enumerative Combinatorics}, Volume 2, Cambridge
University Press, 1999.
\bibitem{fs} P.~Flajolet and R.~Sedgewick, \emph{Analytic Combinatorics}, Cambridge
University Press, 2009.
\bibitem{hw} G.~H.~Hardy and E.~M.~Wright, \emph{An Introduction to the Theory of
Numbers}, 6th ed., Oxford University Press, 2008.
\end{thebibliography}

\end{document}
